% Part 1 导言：表示与渲染
\vspace{1em}
\begin{center}
\textcolor{techgray}{\small 本部分建立3DGS的技术基础——从高斯基元表示、渲染管线到质量优化，是理解全书后续内容的基石。}
\end{center}
\vspace{0.5em}

\noindent\textbf{本部分包含：}
\begin{itemize}[leftmargin=2em,itemsep=1pt]
\item[\textcolor{techblue}{\textbullet}] \textbf{第1章}（3DGS基础）— 高斯基元表示、投影渲染、训练优化、密度控制的技术原理
\item[\textcolor{techblue}{\textbullet}] \textbf{第2章}（渲染与质量）— 抗锯齿、密度控制、训练策略、硬件加速、超分辨率等11个方向
\item[\textcolor{techblue}{\textbullet}] \textbf{第3章}（基元设计）— 纹理高斯、替代核函数、多边形基元、频谱扩展
\end{itemize}

\noindent\textbf{阅读建议：}第1章是全书入口，建议所有读者首先阅读。第2章和第3章可根据研究方向选读——追求渲染质量的读者关注第2章，研究基元改进的读者关注第3章。

\noindent\textbf{与其他部分的联系：}第3章的纹理高斯为第5章（材质分解）奠定基础；第2章的密度控制策略在第11章（模型压缩）中被广泛采用；第2章的排序优化与第14章（大规模重建）共享底层管线。

\vspace{1em}
